\documentclass[11pt]{article}
\usepackage[margin=1in]{geometry}
\usepackage[T1]{fontenc}
\usepackage{lmodern,amsmath,amssymb,booktabs,graphicx}
\usepackage[colorlinks=true,linkcolor=blue,urlcolor=blue]{hyperref}
\setlength{\parindent}{0pt}\setlength{\parskip}{5pt}
\title{RoCE-K: direct calibration of aggregate errors}
\author{}\date{2 October 2026, second experiment update}
\begin{document}\maketitle

\textbf{Scope.} This paired study follows the residual-center intuition
$\tau=d+r_1-r_0$: evaluate compatible residual centers rather than classify
every source. The change is to error calibration, not to aggregation weights.
There are 18 bounded oracle-score settings with 1,000 repetitions each and
12 newly fitted four-stratum settings with 300 each: 21,600 setting-specific
repetitions. Every site contains 1,000 patients. Seeds are new relative to the
preceding report; methods within this study share data. The target audit
reuses the same 3,600 fitted datasets and is not counted as extra repetitions.
All historical outputs, failures and fallback cases remain available.

\section{What changed}

A patient-level, zero-diagonal quadratic representation of the unbiased
source dispersion estimator gives an exponential lower-tail bound. The
result, with its Gaussian MGF comparison proof, is stated in the updated
main design. It uses deterministic score ranges and does not require
simultaneous confidence bounds for every source variance. The pooled mean
uses empirical Bernstein for independent, nonidentically distributed patients.
The new concentration result was checked by a 4,096-atom non-Gaussian
enumeration and independent matrix-spectrum calculations. These checks
support the proof and implementation; they are not a high-dimensional
causal inference theorem.

The fitted diagnostic also controls the signed average nuisance remainder
on \emph{every} subset of the guaranteed size, together with its squared
drift. Source pilot calibration projections, a target outcome-error box,
and joint target probability constraints produce these certificates.
For four strata the required maxima reduce to 16 corners per arm and small
linear programs, checked against an independent LP solver. The algorithm
is linear in source count at fixed strata count, but exponential in strata
count. It is not yet a sparse high-dimensional procedure.

The Gaussian comparison notation is corrected: noncentral chi-square
inversion applies after subtracting the known variance, not a random sample
variance. Historical valid code paths used the former correctly. The revised
API explicitly guards that branch.

\section{Paired bounded-score ablation}

The three source modules are Fourier-only, dispersion-only, and their hybrid.
For every method the total failure allowance is .05. The binary study assigns
.025 to the target, reserves .005 for variance certificates, and assigns
.01 to each arm. The hybrid divides each arm's allowance between Fourier
and dispersion; a single-component method uses the entire arm allowance.
The direct-dispersion procedure reserves the variance allowance even though
it does not need it. Thus its advantage is not due to reclaiming that budget.
Sources have 250 pilot and 750 evaluation patients; the target has 1,000
evaluation patients. Common outcome and weighting functions are known.

The table reports mean interval length relative to the ordinary target-normal
interval at $K=2048$. Weak bias affects one quarter of treated source arms,
with magnitude $2K^{-1/4}$ times a source residual standard error. Control
sources are valid. Both arms are only assumed to have at least $0.75K$
valid sources.

\begin{center}\small
\begin{tabular}{llrrr}\toprule
Target heterogeneity & Source pattern & Previous hybrid & Direct hybrid & Direct dispersion\\\midrule
0 & All valid & .821 & .309 & .297\\
0 & Weak bias & .823 & .313 & .301\\
.3 & Weak bias & 1.736 & 1.079 & 1.066\\
.3 & Strong bias & 1.854 & 1.35 & 1.339\\\bottomrule
\end{tabular}
\end{center}
The exact machine-generated metrics accompany this note; rounded table values
are descriptive. Direct calibration materially improves the bounded-score
result. Shared target prediction variability still limits its usefulness.
Fourier does not show a net length advantage under the fully protected
bounded-score calibrations tested here. This does not imply that Fourier
cannot help in other distributions or under tighter error control.

\begin{figure}[ht]\centering
\includegraphics[width=\linewidth]{bounded_comparison.pdf}
\caption{Weak-bias comparison on identical data. Each method uses the same
.05 total error budget. The dotted line denotes ordinary target-normal length.}
\end{figure}

\section{Fresh fitting: better protection, but no established borrowing gain}

The fitted diagnostic uses four covariate strata. Common outcome and target
propensity fits use half-count smoothing on 500 target training patients;
source joint weights use 500 source training patients. Sources retain 250
pilot and 250 evaluation patients; the target has 500 evaluation patients.
The target truth is .10. This is the same finite-stratum DGP as the preceding
report, with new seeds, not a production sparse-RoCE simulation.

At $K=128$, weak bias and target heterogeneity .3, the average source budget
falls from .285/.263 in the treated/control arms to .111/.097. The oracle
worst-subset signed drifts are only .0034/.0032. Hence the new certificate
is substantially tighter but remains far from the actual drift scale.

\begin{center}\small
\begin{tabular}{lr}\toprule
Procedure, same fitted setting & Mean interval length\\\midrule
Previous hybrid & 1.1101\\
Direct dispersion, previous nuisance budgets & 1.1101\\
Direct hybrid, aggregate nuisance budgets & .8817\\
Direct dispersion, aggregate nuisance budgets & .8722\\
Target evaluation half, matched directional calibration & .9244\\
Target evaluation half, full target-only error allowance & .8781\\
Full target two-fold estimator, equally improved certificate & .8312\\\bottomrule
\end{tabular}
\end{center}

The last comparison changes the substantive conclusion. A shorter interval
than the old target certificate is not sufficient evidence of borrowing
benefit. Applying the same joint-probability nuisance improvement to both
target folds yields a target-only interval still shorter than the proposed
source procedures in this setting. It uses .0025 nuisance and .0225 noise
allowance per fold; the two dependent folds are combined by a union bound.
No source observations enter that target estimator or certificate.

\begin{figure}[ht]\centering
\includegraphics[width=\linewidth]{fitted_comparison.pdf}
\caption{Weak-bias fitted diagnostic with target heterogeneity .3. The left
panel uses the equally improved full-target certificate as denominator.
The right panel distinguishes computable and oracle nuisance quantities.}
\end{figure}

\section{Coverage checks and the remaining decision}

Every source procedure has observed coverage 1 in these settings; the methods
remain conservative. With no misses, a two-sided 95\% binomial lower bound
is about .9963 for 1,000 repetitions and .9878 for 300. This is not evidence
that achieved coverage is exactly .95. All 3,600 fitted repetitions satisfy
the audited mean/squared nuisance and target drift certificates, allowing
$10^{-12}$ numerical tolerance only in that diagnostic.

The original fitted diagnostic compared the theoretically zero prediction
coordinate drift with its floating-point residual (about $10^{-16}$), causing
spurious certificate-failure flags. The original records are preserved and
the separate paired audit supersedes those flags. No interval, point estimate,
or coverage result was altered by this correction.

The next priority is to control the full calibration remainder more sharply
and extend its certificate beyond a small fixed stratum space. Direct
bounded-score dispersion is the simpler useful comparison in this study;
Fourier remains an explicit ablation. Neither large-$K$ coverage reductions
nor successful ideal-score length results establish useful fitted-score
lengths. The target floor, nuisance remainder, and honest target-only
comparison must remain part of the assessment.
\end{document}
